\scanpage{123}
sans restriction à la car. $\ne2$.

\uline{2} \quad On va définir une réflexion canonique dans $P$, ayant $H$ comme hyperplan de pts fixes, et invariant $Q$. Une \inserted{\unsure{pareille}} réflexion \deleted{\unsure{est associée} à un autom} ayant $H$ comme hyperplan fixe \deleted{s'écrit sous la forme} \inserted{induit $E$, donc est \unsure{associé|associée} à un autom. linéaire} $\unsure{\sigma|u}$ de la structure d'extension \deleted{$E$, \unsure{à la}} $T$ de $\mathcal O_S$ par $V$, \deleted{et} induisant l'identité sur $\mathcal O_S$, et une homothétie \inserted{$\lambda\,\mathrm{id}_V$} sur $V$, ou encore (quitte à diviser par $\lambda$\deleted{\unsure{\varphi}}) par un autom de la structure d'ext. induisant l'identité sur $V$, donc de la forme
\[
\sigma=\mathrm{id}+v
\]
où $v$ est nul sur $V$, donc de la forme
\[
v(x)=\pi(x).c\qquad\text{i.e.}\quad v=\pi\otimes c
\]
où $\pi\in\Gamma(S,\check T)$ est l'hom $T\longrightarrow\mathcal O_S$ définissant la structure d'extension, $c$ une section quelc. de $T$. Écrivons que \deleted{\unsure{Q}} $u$ invarie $q$ mod homothétie, on trouve
\[
\begin{aligned}
q\circ\sigma(x)&=q(x+vx)=q(x+\pi(x).c)\\
&=q(x)+\pi(x)^2q(c)+\deleted{\unsure{q}}\pi(x)\varphi(x,c)
\end{aligned}
\]
On veut trouver \deleted{\unsure{x}}$\lambda\in\Gamma(S,G_m)$ tel que
\[
\deleted{qu(x)=\lambda q(x)}\quad q\sigma(x)=\lambda q(x)\qquad\forall x\in\Gamma(-,T)
\]
% archive sheet number: 122
